% sections/05-code-modularization.tex — Code Modularization

\section{Code Modularization}

% ── Requirement ──
\subsection{Requirement}

Your code must be properly modularized. The following are \textbf{NOT} acceptable:

\begin{itemize}
  \item A single file with thousands of lines of code
  \item All code placed in \filename{main.cpp}
  \item Duplicated code across files
  \item Tightly coupled components that cannot be tested independently
\end{itemize}

% ── Project Structure ──
\subsection{Project Structure}

Your project should follow this directory structure:

\begin{lstlisting}[style=tree]
project/
|-- include/                    Header files (.h/.hpp)
|   |-- data_structures/
|   |-- core/
|   +-- utils/
|-- src/                        Source files (.cpp)
|   |-- data_structures/
|   |-- core/
|   |-- utils/
|   +-- main.cpp
|-- tests/                      Test files
|   |-- unit/
|   |-- integration/
|   |-- edge_cases/
|   |-- performance/
|   +-- fixtures/
|-- data/                       Input/test data
|   |-- input_*.json
|   |-- edge_case_*.xml
|   |-- test_*.yaml
|   |-- expected_*.json
|   +-- benchmark_*.bin
|-- scripts/
|   |-- generate_test_data.py
|   +-- generate_binary.cpp
|-- docs/
|   |-- uml/
|   +-- Doxyfile
|-- frontend/                   Frontend implementation
|-- libs/                       External header-only libraries
|-- Makefile                    (or CMakeLists.txt)
|-- CMakeLists.txt              (optional, alternative to Makefile)
|-- project.code-workspace      VSCode workspace (optional)
|-- .clang-format
+-- README.md
\end{lstlisting}

% ── Modularization Guidelines ──
\subsection{Modularization Guidelines}

\subsubsection{Single Responsibility Principle}

Each file and class should have a single, well-defined purpose. For example:

\begin{itemize}
  \item \filename{LinkedList.h} / \filename{LinkedList.cpp} --- linked list data structure
  \item \filename{FileParser.h} / \filename{FileParser.cpp} --- file parsing logic
  \item \filename{Validator.h} / \filename{Validator.cpp} --- input validation
\end{itemize}

\subsubsection{Utility Functions}

Group reusable helper functions into utility modules:

\begin{lstlisting}[style=cpp]
// utils/StringUtils.h
#pragma once
#include <string>
#include <vector>

namespace utils {

std::string trim(const std::string& str);
std::string toLower(const std::string& str);
std::vector<std::string> split(const std::string& str, char delimiter);
bool startsWith(const std::string& str, const std::string& prefix);

} // namespace utils
\end{lstlisting}

\textit{\textbf{Note:} Utility functions may return STL types (e.g., \texttt{std::vector}) for transient use. This does not violate the STL restriction, which applies to your \textbf{primary data structures} --- the ones central to your solution.}

\subsubsection{Common Patterns}

Provide a centralized logging utility:

\begin{lstlisting}[style=cpp]
// utils/Logger.h
#pragma once
#include <string>

enum class LogLevel { DEBUG, INFO, WARNING, ERROR };

void log(LogLevel level, const std::string& message);
void setLogLevel(LogLevel level);
\end{lstlisting}

\subsubsection{Configuration}

Centralize application configuration:

\begin{lstlisting}[style=cpp]
// core/Config.h
#pragma once
#include <string>

struct Config {
    std::string dataDirectory;
    int maxIterations;
    bool verboseMode;

    static Config loadFromFile(const std::string& path);
};
\end{lstlisting}

% ── File Size Guidelines ──
\subsection{File Size Guidelines}

\begin{center}
\begin{tabular}{l l}
\hline
\rowcolor{tableheader} \textbf{File Type} & \textbf{Recommended Max Lines} \\
\hline
Header files (\filename{.h}) & 200--300 lines \\
Source files (\filename{.cpp}) & 300--500 lines \\
\filename{main.cpp} & 50--100 lines \\
Test files & 200--400 lines \\
\hline
\end{tabular}
\end{center}

If a file exceeds these limits, consider splitting it into smaller, focused modules.

% ── SOLID Design Principles ──
\subsection{SOLID Design Principles (Recommended)}

\textit{The following principles are recommended for better code quality but not strictly required. Focus on basic modularization first.}

\subsubsection{Single Responsibility Principle (SRP)}

Each class or module should have only one reason to change. A class that handles both file parsing and data validation is doing too much---split it.

\subsubsection{Open/Closed Principle (OCP)}

Software entities should be open for extension but closed for modification. Use inheritance and polymorphism to add new behavior without changing existing code.

\subsubsection{Liskov Substitution Principle (LSP)}

Derived classes must be substitutable for their base classes without altering program correctness.

\begin{lstlisting}[style=cpp]
class Shape {
public:
    virtual ~Shape() = default;
    virtual double area() const = 0;
    virtual double perimeter() const = 0;
    virtual void draw() const = 0;
};

class Circle : public Shape {
    double radius;
public:
    explicit Circle(double r) : radius(r) {}
    double area() const override;
    double perimeter() const override;
    void draw() const override;
};

class Rectangle : public Shape {
    double width, height;
public:
    Rectangle(double w, double h) : width(w), height(h) {}
    double area() const override;
    double perimeter() const override;
    void draw() const override;
};

class Triangle : public Shape {
    double a, b, c;
public:
    Triangle(double a, double b, double c) : a(a), b(b), c(c) {}
    double area() const override;
    double perimeter() const override;
    void draw() const override;
};

// Works with ANY Shape --- LSP in action
void printShapeInfo(const Shape& shape) {
    std::cout << "Area: " << shape.area() << "\n";
    std::cout << "Perimeter: " << shape.perimeter() << "\n";
    shape.draw();
}
\end{lstlisting}

\subsubsection{Interface Segregation Principle (ISP)}

Prefer small, focused interfaces over large monolithic ones. Clients should not be forced to depend on methods they do not use.

\subsubsection{Dependency Inversion Principle (DIP)}

High-level modules should not depend on low-level modules. Both should depend on abstractions. Abstractions should not depend on details---details should depend on abstractions.

% ── Design Patterns (Optional) ──
\subsection{Design Patterns (Optional --- for bonus points)}

\textit{Design patterns can improve code architecture. Only use them if they genuinely fit your problem.}

\subsubsection{Factory Pattern}

\begin{lstlisting}[style=cpp]
class VehicleFactory {
public:
    enum class Type { CAR, MOTORCYCLE, TRUCK, BICYCLE };

    static std::unique_ptr<Vehicle> create(Type type) {
        switch (type) {
            case Type::CAR:        return std::make_unique<Car>();
            case Type::MOTORCYCLE: return std::make_unique<Motorcycle>();
            case Type::TRUCK:      return std::make_unique<Truck>();
            case Type::BICYCLE:    return std::make_unique<Bicycle>();
            default:               return nullptr;
        }
    }

    static std::unique_ptr<Vehicle> create(const std::string& type) {
        if (type == "car")        return create(Type::CAR);
        if (type == "motorcycle") return create(Type::MOTORCYCLE);
        if (type == "truck")      return create(Type::TRUCK);
        if (type == "bicycle")    return create(Type::BICYCLE);
        return nullptr;
    }
};
\end{lstlisting}

\subsubsection{Strategy Pattern}

\begin{lstlisting}[style=cpp]
class IPaymentStrategy {
public:
    virtual ~IPaymentStrategy() = default;
    virtual void pay(double amount) = 0;
    virtual std::string getName() const = 0;
};

class CreditCardPayment : public IPaymentStrategy {
public:
    void pay(double amount) override;
    std::string getName() const override { return "Credit Card"; }
};

class PayPalPayment : public IPaymentStrategy {
public:
    void pay(double amount) override;
    std::string getName() const override { return "PayPal"; }
};

class CryptoPayment : public IPaymentStrategy {
public:
    void pay(double amount) override;
    std::string getName() const override { return "Crypto"; }
};

class ShoppingCart {
    std::unique_ptr<IPaymentStrategy> paymentStrategy;
public:
    void setPaymentStrategy(std::unique_ptr<IPaymentStrategy> strategy) {
        paymentStrategy = std::move(strategy);
    }
    void checkout(double total) {
        paymentStrategy->pay(total);
    }
};
\end{lstlisting}

\subsubsection{Observer Pattern}

\begin{lstlisting}[style=cpp]
class IObserver {
public:
    virtual ~IObserver() = default;
    virtual void onEvent(const std::string& event) = 0;
};

class EventManager {
    LinkedList<IObserver*> observers;  // Your custom linked list
public:
    void subscribe(IObserver* observer) {
        observers.add(observer);
    }
    void unsubscribe(IObserver* observer) {
        observers.remove(observer);
    }
    void notify(const std::string& event) {
        // Iterate over observers and call onEvent
        for (auto it = observers.begin(); it != observers.end(); ++it) {
            (*it)->onEvent(event);
        }
    }
};
\end{lstlisting}

% ── Module Organization by Feature ──
\subsection{Module Organization by Feature}

For larger projects, organize source code by feature or domain rather than by file type:

\begin{lstlisting}[style=tree]
src/
|-- inventory/
|   |-- Inventory.h
|   |-- Inventory.cpp
|   |-- Item.h
|   +-- ItemFactory.h
|-- entities/
|   |-- Entity.h
|   |-- Player.h
|   +-- Enemy.h
+-- combat/
    |-- CombatManager.h
    +-- strategies/
        |-- IAttackStrategy.h
        |-- MeleeAttack.h
        +-- RangedAttack.h
\end{lstlisting}

% ── Dependency Management ──
\subsection{Dependency Management}

\begin{itemize}
  \item \textbf{Minimize dependencies} between modules. Each module should depend on as few others as possible.
  \item \textbf{Use forward declarations} where possible to reduce include chains.
  \item \textbf{Avoid circular dependencies.} If \filename{A.h} includes \filename{B.h} and \filename{B.h} includes \filename{A.h}, restructure your code.
  \item \textbf{Document dependencies} in header comments so other developers can understand the dependency graph.
\end{itemize}
